\documentclass{article}
\usepackage[utf8]{inputenc}
\usepackage{xcolor}
\usepackage{listings}
\usepackage{tcolorbox}
\usepackage{blindtext}
\usepackage{textcomp}
\usepackage{amsmath}
\usepackage{amsfonts}
\usepackage{amssymb}
\usepackage{bm}
\usepackage{physics}
\usepackage{algorithm2e}
\usepackage{tcolorbox}
\usepackage{hyperref}
\usepackage{tabularx}
\usepackage{appendix}
\usepackage{bbm}

\usepackage{color}   %May be necessary if you want to color links
\usepackage{hyperref}
\hypersetup{
    colorlinks=true, %set true if you want colored links
    linktoc=all,     %set to all if you want both sections and subsections linked
    linkcolor=blue,  %choose some color if you want links to stand out
}

\definecolor{light-gray}{gray}{0.95}
\definecolor{codegreen}{rgb}{0,0.6,0}
\definecolor{codegray}{rgb}{0.5,0.5,0.5}
\definecolor{codepurple}{rgb}{0.58,0,0.82}
\definecolor{backcolour}{rgb}{0.95,0.95,0.92}
\lstdefinestyle{mystyle}{
    backgroundcolor=\color{backcolour},   
    commentstyle=\color{codegreen},
    keywordstyle=\color{magenta},
    numberstyle=\tiny\color{codegray},
    stringstyle=\color{codepurple},
    basicstyle=\ttfamily\footnotesize,
    breakatwhitespace=false,         
    breaklines=true,                 
    captionpos=b,                    
    keepspaces=true,                 
    numbers=left,                    
    numbersep=5pt,                  
    showspaces=false,                
    showstringspaces=false,
    showtabs=false,                  
    tabsize=2
}
\lstset{style=mystyle,literate={~} {$\sim$}{1}}
\newcommand{\code}[1]{\colorbox{light-gray}{\texttt{#1}}}

\newcommand{\R}{\mathbb{R}}
\newcommand{\N}{\mathbb{N}}
\newcommand{\Z}{\mathbb{Z}}
\newcommand{\Q}{\mathbb{Q}}
\newcommand{\C}{\mathbb{C}}
\newcommand{\proof}{\textbf{Proof: }}
\newcommand{\theorem}{\textbf{Theorem: }}

\renewcommand{\var}[1]{\mathrm{Var}\ensuremath{\left[ #1 \right]}}

\let\StandardTheFigure\thefigure
\let\vec\mathbf
\providecommand{\cov}[1]{\mathbf{Cov}\ensuremath{\left( #1 \right)}}
\providecommand{\pr}[1]{\ensuremath{\Pr\left(#1\right)}}
\providecommand{\brak}[1]{\ensuremath{\left(#1\right)}}
\providecommand{\cbrak}[1]{\ensuremath{\left\{#1\right\}}}
\providecommand{\sbrak}[1]{\ensuremath{\left[#1\right]}}
\providecommand{\mean}[1]{\mathbb{E}\ensuremath{\left[ #1 \right]}}
%\providecommand{\var}[1]{\mathbb{Var}\ensuremath{\left[ #1 \right]}}
\providecommand{\der}[1]{\mathrm{d} #1}
\providecommand{\gauss}[2]{\mathcal{N}\ensuremath{\left(#1,#2\right)}}
\providecommand{\mbf}{\mathbf}
\providecommand{\abs}[1]{\ensuremath{\left\vert#1\right\vert}}
\providecommand{\norm}[1]{\ensuremath{\left\lVert#1\right\rVert}}
\providecommand{\z}[1]{{\mathcal{Z}}\{#1\}}
\providecommand{\ztrans}{\overset{\mathcal{Z}}{ \rightleftharpoons}}
\providecommand{\supp}[1]{\mathrm{supp}\ensuremath{\left(#1\right)}}
\providecommand{\cond}[2]{\ensuremath{p(#1\,|\,#2)}}

\newcommand{\myvec}[1]{\ensuremath{\begin{pmatrix}#1\end{pmatrix}}}
\newcommand{\mydet}[1]{\ensuremath{\begin{vmatrix}#1\end{vmatrix}}}
\newcommand{\define}{\stackrel{\triangle}{=}}
\newcommand{\solution}{\textbf{\underline{Solution :}}}
\DeclareMathOperator*{\argmin}{arg\,min}
\DeclareMathOperator*{\argmax}{arg\,max}

\title{\textbf{Mphasis Optimization Formulation}}

\author{
Abhay Shankar K | CS21BTECH11001 \\
Anirudh Srinivasan | CS20BTECH11059 \\
Gautam Singh | CS21BTECH11018
}
\date{6th December 2023}

\begin{document}
\maketitle
\newpage

%\tableofcontents

\begin{flushleft}
%-----------------------------------------------
%\begin{abstract}
%\end{abstract}
%\newpage
\section{\textbf{Rules}}
% Rule Set for PNR Scoring (PNR Individual)
\subsection{Rule Set for PNR Scoring (PNR Individual)}
\begin{enumerate}
    \item \textbf{SSR (Special Service Request):}
    \begin{itemize}
        \item \textbf{Name:} SSR
        \item \textbf{Condition:} PNR.SSR = SSR
        \item \textbf{Score:} 200
    \end{itemize}
    
    \item \textbf{Cabin:}
    \begin{itemize}
        \item \textbf{Name:} Cabin
        \item \textbf{Conditions:} PNR.Cabin = F, J, Y
        \item \textbf{Score:} 1500 - 2000 (Dynamic scoring)
    \end{itemize}

    % Add other rules for Class, Connection, Paid Service, Booking Type, Number of Pax, Loyalty
\end{enumerate}

% Rule Set for Flight Scoring
\subsection{Rule Set for Flight Scoring}
\begin{enumerate}
    \item \textbf{Arrival Delay:}
    \begin{itemize}
        \item \textbf{Name:} Arrival Delay
        \item \textbf{Conditions:} Flight.proposed Arrival Delay $\leq$ 6/12/24/48 hours
        \item \textbf{Score:} 70/50/40/30 (Sample scoring)
    \end{itemize}

    % Add other rules for Equipment, City Pairs, STD of Proposed Flight, Stopover
\end{enumerate}

% Quality Score Grading for Flight Selection
\subsection{Quality Score Grading for Flight Selection}
% Add grading table here

% Rules for Alternate Flight Selection and Connections
\subsection{Rules for Alternate Flight Selection and Connections}
% Add rules for Propose Flight, Segment, Grouping Rules

% Downgrade and Upgrade Rules
\subsection{Downgrade and Upgrade Rules}
% Add rules for Class Mapping, Downgrade Class Mapping, Upgrade Class Mapping

% Cabin Class Mapping
\subsection{Cabin Class Mapping}
% Add cabin class mapping table

% Solution Ranking Factors
\subsection{Solution Ranking Factors}
% Add factors and notes on score values

% Journey-Based Solution Scoring Mechanism
\subsection{Journey-Based Solution Scoring Mechanism}
% Add scoring mechanism details

% \section{Optimization}
% \subsection{Objective Function}
% Minimize the overall cost and impact of passenger re-accommodation, considering various factors such as time, service requests, loyalty, class preferences, and flight characteristics.

% \subsection{Decision Variables}
% \begin{enumerate}
%     \item Binary variables indicating the selection of each alternate flight for re-accommodating passengers.
% \item Binary variables representing the upgrade or downgrade of passenger classes.
% \end{enumerate}


% \subsection{Constraints}
% \begin{enumerate}
%     \item \textbf{SSR Constraints:} Ensure that special service requests (SSRs) are accommodated according to individual scoring rules.

% \item \textbf{Cabin Constraints:} Accommodate passengers in suitable cabins based on dynamic scoring for each cabin class.

% \item \textbf{Class Constraints:} Assign passengers to appropriate classes considering scoring rules for different passenger classes.

% \item \textbf{Connection Constraints:} Manage segments, ensuring the connection time adheres to specified rules.

% \item \textbf{Paid Service Constraints:} Consider the impact of paid services on passenger re-accommodation.

% \item \textbf{Booking-Type Constraints:} Handle different booking types (e.g., corporate, group, military) based on specific scoring rules.

% \item \textbf{Number of Pax Constraints:} Consider the number of passengers in a PNR and ensure proper accommodation based on scoring.

% \item \textbf{Loyalty Constraints:} Accommodate passengers based on loyalty levels, adhering to loyalty-specific scoring rules.

% \item \textbf{Arrival Delay Constraints:} Manage alternate flight proposals based on arrival delay conditions.

% \item \textbf{Equipment Constraints:} Ensure consistency in equipment assignment for proposed flights.

% \item \textbf{City Pairs Constraints:} Consider city pairs for proposed flights, adhering to specified conditions.

% \item \textbf{STD of Proposed Flight Constraints:} Manage proposed flight STD based on original flight STD, considering specified conditions.

% \item \textbf{Stopover Constraints:} Account for stopovers in the decision-making process.

% \item \textbf{Flight Quality Score Grading Constraints:} Classify flights into quality score grades based on the total score obtained.

% \item \textbf{Alternate Flight Selection and Connection Constraints:} Apply conditions for proposing or rejecting alternate flights based on specific rules.

% \item \textbf{Downgrade and Upgrade Constraints:} Implement rules for class mapping in case of upgrades or downgrades.

% \item \textbf{Cabin Class Mapping Constraints:} Utilize cabin class mapping rules for identifying suitable cabins for each class.

% % Add these to objective function formulation
% %\item \textbf{Solution Ranking Constraints:} Evaluate solution ranking based on defined factors and scoring mechanisms.

% %\item \textbf{Journey-Based Solution Scoring Constraints:} Consider the number of solution combinations and score each solution file accordingly.
% \end{enumerate}

% % \subsection{Constrained Quadratic Model (CQM)}

% % Let's use generic symbols:
% % \begin{enumerate}
% %   \item \( \text{Cost}_{\text{Flight}} \)
% %   \item \( \text{Cost}_{\text{Passenger}} \)
% %   \item \( \text{Score}_{\text{SSR}} \)
% %   \item \( \text{Score}_{\text{Cabin}} \)
% %   \item \( \text{Score}_{\text{Class}} \)
% %   \item \( \text{Score}_{\text{Connection}} \)
% %   \item \( \text{Score}_{\text{PaidService}} \)
% %   \item \( \text{Score}_{\text{BookingType}} \)
% %   \item \( \text{Score}_{\text{NumberOfPax}} \)
% %   \item \( \text{Score}_{\text{Loyalty}} \)
% %   \item \( \text{Score}_{\text{ArrivalDelay}} \)
% %   \item \( \text{Score}_{\text{Equipment}} \)
% %   \item \( \text{Score}_{\text{CityPair}} \)
% %   \item \( \text{Score}_{\text{STD}} \)
% %   \item \( \text{Score}_{\text{Stopover}} \)
% %   \item \( \text{TotalScore}_{\text{Flight}} \)
% %   \item \( \text{Score}_{\text{SolutionRanking}} \)
% %   \item \( \text{Score}_{\text{JourneySolution}} \)
% % \end{enumerate}

% % \subsubsection{Objective Function}
% % \[
% % \text{Minimize} \quad \sum_{\text{flights}} \text{Cost}_{\text{Flight}} \times \text{BinaryVariable}_{\text{Flight}} + \sum_{\text{passengers}} \text{Cost}_{\text{Passenger}} \times \text{BinaryVariable}_{\text{Passenger}}
% % \]

% % % Single Objective Function - Concise
% % % We'll have to look at conflicting objectives:
% % % If yes: Multi-objective Opt

% % \subsubsection{Decision Variables}
% % \begin{enumerate}
% %     \item \(\text{BinaryVariable}_{\text{Passenger, Flight}}\) represents the selection of each alternate flight for re-accommodating passengers.
% % \item \(\text{BinaryVariable}_{\text{Passenger, Class}}\) represents the upgrade or downgrade of passenger classes.
% % \end{enumerate}

% % %

% % \subsubsection{Constraints}

% % \begin{align*}
% % & \sum_{\text{SSR types}} \text{Score}_{\text{SSR}} \times \text{BinaryVariable}_{\text{SSR type}} \geq \text{Threshold}_{\text{SSR}} \\
% % & \sum_{\text{Cabins}} \text{Score}_{\text{Cabin}} \times \text{BinaryVariable}_{\text{Cabin}} \geq \text{Threshold}_{\text{Cabin}} \\
% % & \sum_{\text{Classes}} \text{Score}_{\text{Class}} \times \text{BinaryVariable}_{\text{Class}} \geq \text{Threshold}_{\text{Class}} \\
% % & \text{Score}_{\text{Connection}} \times \text{BinaryVariable}_{\text{Connection}} \geq \text{Threshold}_{\text{Connection}} \\
% % & \text{Score}_{\text{PaidService}} \times \text{BinaryVariable}_{\text{PaidService}} \geq \text{Threshold}_{\text{PaidService}} \\
% % & \text{Score}_{\text{BookingType}} \times \text{BinaryVariable}_{\text{BookingType}} \geq \text{Threshold}_{\text{BookingType}} \\
% % & \text{Score}_{\text{NumberOfPax}} \times \text{BinaryVariable}_{\text{NumberOfPax}} \geq \text{Threshold}_{\text{NumberOfPax}} \\
% % & \text{Score}_{\text{Loyalty}} \times \text{BinaryVariable}_{\text{Loyalty}} \geq \text{Threshold}_{\text{Loyalty}} \\
% % & \sum_{\text{Arrival Delay types}} \text{Score}_{\text{ArrivalDelay}} \times \text{BinaryVariable}_{\text{ArrivalDelay type}} \geq \text{Threshold}_{\text{ArrivalDelay}} \\
% % & \text{Score}_{\text{Equipment}} \times \text{BinaryVariable}_{\text{Equipment}} \geq \text{Threshold}_{\text{Equipment}} \\
% % & \sum_{\text{City Pair types}} \text{Score}_{\text{CityPair}} \times \text{BinaryVariable}_{\text{CityPair type}} \geq \text{Threshold}_{\text{CityPair}} \\
% % & \sum_{\text{STD types}} \text{Score}_{\text{STD}} \times \text{BinaryVariable}_{\text{STD type}} \geq \text{Threshold}_{\text{STD}} \\
% % & \text{Score}_{\text{Stopover}} \times \text{BinaryVariable}_{\text{Stopover}} \geq \text{Threshold}_{\text{Stopover}} \\
% % & \text{TotalScore}_{\text{Flight}} \geq \text{Threshold}_{\text{TotalScore}} \\
% % & \text{BinaryVariable}_{\text{ProposeFlight}} = 0 \quad \text{if} \quad \text{Condition}_{\text{ProposeFlight}} \quad \text{is satisfied} \\
% % & \text{BinaryVariable}_{\text{segment}} = 0 \quad \text{if} \quad \text{Conditions}_{\text{segment}} \quad \text{are satisfied} \\
% % & \text{BinaryVariable}_{\text{ClassMapping}} = 1 \quad \text{if} \quad \text{Conditions}_{\text{ClassMapping}} \quad \text{are satisfied} \\
% % & \sum_{\text{Cabin classes}} \text{BinaryVariable}_{\text{CabinClassMapping}} \geq \text{Threshold}_{\text{CabinClassMapping}} \\
% % & \text{Score}_{\text{SolutionRanking}} \times \text{BinaryVariable}_{\text{SolutionRanking}} \geq \text{Threshold}_{\text{SolutionRanking}} \\
% % & \text{Score}_{\text{JourneySolution}} \times \text{BinaryVariable}_{\text{JourneySolution}} \geq \text{Threshold}_{\text{JourneySolution}} \\
% % \end{align*}

% % One can consider setting the thresholds to all zero as a good indicator.

% \subsection{Discrete Quadratic Model (DQM) Formulation I}
% \subsubsection{Notation}

% \begin{table}[!ht]
%     \centering
%     \begin{tabularx}{\textwidth}{|X|X|X|}
%         \hline
%         \textbf{Variable} & \textbf{Definition} & \textbf{Comments} \\
%         \hline
%         \(f_{ij}\) & Object for flight \(j\) wrt flight \(i\) & Constant \\
%         \hline
%         \(s_i\) & Object for PNR \(i\) & Constant \\
%         \hline
%         \(g_i\) & \(s_i\).flight\_number & Constant \\
%         \hline
%     \end{tabularx}
% \end{table}

% % Binary variables for user to toggle business rules
% % 1. some handled by preprocessing.
% % 2. opt: add some for user convenience (cabin/class upgrade/downgrade), given
% %    by preprocessing.
% % 3. binary variable for PNR to flight? b[i][j]: 1 if PNR i gets flight j.
% % 4. Flight selection -> constraints
% % 5. Implement solution ranking as constraints in the lp

% % b[i][j] sum over j to 1.
% % b[i][j] sum over i to capacity[j].
% % For a ranking list, sum over j is top k.
% % What about the sum over i?

% \subsubsection{Decision Variables}

% We define our decision variables to be \(b_{ij}\), such that

% \begin{equation}
%     b_{ij} =
%     \begin{cases}
%         1 & \textrm{passenger \(i\) is allocated to flight \(j\)} \\
%         0 & \textrm{otherwise}
%     \end{cases}.
%     \label{eq:bij-def}
% \end{equation}

% % We also define additional decision variables \(r_i\), where

% % \begin{equation}
% %     r_i =
% %     \begin{cases}
% %         1 & \textrm{business rule \(i\) is toggled on} \\
% %         0 & \textrm{otherwise}
% %     \end{cases}.
% %     \label{eq:ri-def}
% % \end{equation}

% \subsubsection{Objective Function}

% Considering this problem to be a maximal matching problem, the objective
% function for this problem is

% % Objective Function
% % - Score for passengers allocated to required flights.
% % - Score for overbooking, using delay as the cost factor.
% % - Score for journey-based solution mechanism: Try to map one flight to
% %   another, but not one flight to multiple other flights. Penalize out-degree of
% %   the bipartite passenger -> flight directed graph?
% %   * Feels like cardinality of the union of j where the flight number of all
% %     PNRs i is same.

% \begin{equation}
%     E\brak{\vec{b}} = \sum_{i,j}b_{ij}\brak{s_i.\text{score} + f_{g_ij}.\text{score} - \lambda\sbrak{f_{g_ij}.\text{eta} - g_i.\text{eta}}_+} - \sum_{j}\mathbbm{1}_{\cbrak{\sum_{i}b_{ij} > 0}}
%     \label{eq:obj-def}
% \end{equation}

% \subsubsection{Constraints}

% The constraints for this problem are

% \begin{align}
%     % Each passenger boards at most one flight.
%     \sum_{j}b_{ij} &\le 1\ \forall\ i \\
%     % Should not exceed flight capacity.
%     \sum_{i}b_{ij}\brak{s_{i}.\text{cap}} &\le \text{cap}_j\ \forall\ j \\
%     % Alternate flight ETD should not cross a threshold.
%     \mathbbm{1}_{\cbrak{\sum_{i}b_{ij} > 0}}\brak{f_{g_ij}.\text{etd} - g_i.\text{std}} &\le t_d\ \forall\ i,j \\
%     % Segment time gap must be between thresholds.
%     \max_{j : \sum_{i}b_{ij} > 0}g_i.\text{next}.\text{std} - f_{g_ij}.\text{eta} &\le t_{d2}\ \forall\ i,j \\
%     \min_{j : \sum_{i}b_{ij} > 0}g_i.\text{next}.\text{std} - f_{g_ij}.\text{eta} &\ge t_{d1}\ \forall\ i,j
% \end{align}

\section{DQM}

% Really should have thought about this before. 

This is not a strict specification. There may be some inconsistencies. Naturally, the version in \texttt{final.ipynb} is the correct one.

We are given a vector of PNR's \texttt{PNR}, Flights \texttt{Flights}, and cancelled flights \texttt{Cancelled}. %\texttt{Ded}.

\subsection{Things we compute}.

\begin{itemize}
    \item Airport-flight directed graph, \(\mathcal{A}\). Vertices are airports, edges are flights. Represent as adjacency matrix. \begin{itemize}
        \item Each edge has a weight, which is the score of the flight. 
        \item For representing different citypair travel, add edges with negative weights between every pair of nodes? This precludes us from using simpler routing algos.
        \item Maximise weight instead of minimising. Might still need DP.
    \end{itemize}

\[A_{ij} = \begin{cases}
    1 & \exists f : \textrm{f.origin = i } \wedge \textrm{ f.destination = j} \\
    0 & \textrm{otherwise}
\end{cases}\]

    \item Flight-flight directed graph, \(\mathcal{F}\). Vertices are flights, edges exist if passenger can get off one and onto another. Represent as adjacency matrix.
    
    This need not be a preprocessing step, but would appear to be most convenient there.


\[F_{ij} = \begin{cases}
    1 & \textrm{Flights[j].std - Flights[i].eta } \in \sbrak{1, 12} \wedge \textrm{Flights[j].origin = Flights[i].destination} \\ % numbers subject to change
    0 & \textrm{otherwise}
\end{cases}\]
    \item PNR vs flight path map, \(\mathcal{P}\). This may be represented as a matrix of PNR vs Flight, where the element is 1 if the flight is in the path for the PNR, and 0 otherwise. \[P_{ij} = \begin{cases}
        \textrm{1} & \textrm{Flight[j] \(\in\) PNR[i].flightpath} \\ % much more concise in code
        0 & \textrm{otherwise}
    \end{cases}\]

    \textbf{Sequence number thing did not work.}

    Note: This is the solution. As such, none of the flights should be in the cancelled list.
    \item A list of impacted PNR indices, \(\mathcal{P'}\).
    
    \[\forall i \in PNR,\ i \in P' \iff i.flightpath \cap Cancelled \neq \phi \]
    \item A flight-flight graph relating arrival times (refer code), \(D\) (D for delta).
\end{itemize}

\subsection{Constraints}.

\begin{itemize}
    \item Flight cancellation enforced: \[\sum_{i, j} \mathbbm{1}_{Flights[j] \in Cancelled} P_{ij} = 0\]
    % \item Source-destination constraints: \begin{align*}
    %     \forall i\ \exists f : f &= \argmax_j\cbrak{P_{ij}.eta} \wedge f.destination = PNR[i].destination \\
    %     \forall i\ \exists f : f &= \argmin_j\cbrak{P_{ij}.eta} \wedge f.origin = PNR[i].origin 
    % \end{align*}
    \item Flight capacity constraint. \[\sum_{i}P_{ij}PNR[i].pax \leq \textrm{Flights[j].capacity}\]
    \item Flight connection constraint: Each PNR's flightpath must be in the flight-flight graph. 
    
    % Every edge proposed must be a part of the flight graph. So, we make a variable for each \emph{pair} of consecutive flights (think 2-window) and check existence. Let \(F'[i]\) be the sequence of 2-windows for the flightpath of the i'th PNR.


    \begin{itemize}
        \item Discrete variable cases \begin{itemize}
            \item Say maximum 5 segments allowed. Each \(P_{ij}\) has 5 variants \(1..=5\) and the 0, and each must appear at most once. \[\forall i \sum_{j} P_{ij}::k \leq 1\ \forall k\]
            \item Segment k cannot appear without segment k - 1. \[\forall i \sum_{j} P_{ij}::k \leq \sum_{j} P_{ij}::k - 1\ \forall k \in 2..5\]
            \item If segment k appears, then flight connection. For every i, x, y, if \(P_{ix}::k\) and \(P_{iy}::k + 1\) then \(F_{xy} = 1\). 
            \[\forall i, k,\ \sum_{x, y} P_{ix}::k \cdot P_{iy}::(k + 1) \cdot (F_{xy} - 1) = 0\]

            This cannot be expressed as a constraint, so it will be added to the objective function we seek to \emph{maximize}, as a term \[\forall i,k,\ \lambda_c \sbrak{\sum_{x, y} P_{ix}::k \cdot P_{iy}::(k + 1) \cdot (F_{xy} - 1)}\]

            We need to provide a high \(\lambda_c\) to reduce likelihood of violation.

            \emph{Feature, not bug}: Sometimes, this may get violated if score is very high. This is good.
        \end{itemize}
    \end{itemize}

    \item Flight delay constraint: 

    \[\forall i \in \mathcal{P'}\ \bigwedge_j D[Flights[j]][PNR[i].last] == 1\]
    \item XXED: When VPNR has XX'ed a flight, should not get reassigned to that flight.
\end{itemize}

\subsection{Objective Function}

\[E\brak{whatever} = \sum_{i} \brak{ PNR[i].score \cdot T_0 - \lambda_1 T_1 + \lambda_2 T_2}\]

where

\begin{align*}
    T_0 &=  \sum_{j} P_{ij} \brak{Flight[j].score - \lambda_0} \\
    T_1 &= \sum_j P_{ij}(Flights[j].eta - PNR[i].sta) \\
    T_2 &= \cdot \sum_{j} \mathbbm{1}_{C} P_{ij} \\
    C &= Flights[j].origin == PNR[i].origin \\\|& Flights[j].destination == PNR[i].destination
\end{align*}

\(\lambda_0\) penalizes number of flights \textit{per PNR}. The value of \(\lambda_0\) should be such that two flights combined can never have a higher score than a single flight.


\section{Transplant}

A Discrete Quadratic Model is an optimisation model where each decision variable may have one or more variants (aka cases) with an objective function of the form 
\[Z = \sum_{i} a_i(x_i) + \sum_{i} \sum_{j} b_{ij} (x_i, x_j) + C\]

where \begin{itemize}
    \item \(x_i\) are the decision variables
    \item \(a_i\) and \(b_{ij}\) are functions that map between the variants of the decision variables.
\end{itemize}

We have used such a model to solve the given problem, employing several intermediate data structures for a clean formulation. These are \begin{itemize}
    \item \(\mathcal{F}\): A directed graph where each node is a flight, and an edge \(a \rightarrow b\)exists if the following conditions hold. \begin{itemize}
        \item The destination of \(a\) must be the source of \(b\)
        \item The departure time of \(b\) is within 1 to 12 hours after the arrival time of \(a\).
    \end{itemize}
    \item \(\mathcal{D}\): Another directed flight graph. An edge \(a \rightarrow b\) exists if The departure time of \(b\) is less than 72 hours after the arrival time of \(a\).
    \item \texttt{PNR}: A list of virtual PNRs. Each PNR in the input is split into multiple contiguous sequences of flights, each of which corresponds into a VPNR.
    \item \texttt{flights}: A list of flights. Read from input.
    \item We also accept a list of cancelled flights from the user.
\end{itemize}

The output of the model is \(\mathcal{P}\), a VPNR-vs-Flight matrix. \(\mathcal{P}_{ij}\) is the sequence number of flight j in the journey of VPNR i, or 0 not applicable.

\end{flushleft}
\end{document}

%%%%%%%%%
% Flight level granularity
% Graphs: (Directed) f1 -> f2 if f2 can replace f1! ()
% - Airport graph, replace time after...?
% - f1 -> ? -> f3 // replacements
% - constraints are capacities (sum of in-degree weights)
% Weights: Time delay
% Each flight has some number of passengers who want to board it.
% Then, list of cancelled flights.
% Find alternatives for each flight.
% 
%%%%%%%%%